\clearpage\pagestyle{otMain}\fontsize{11}{15}\selectfont
\corrigesChapitre{11}{Proportions et évolutions}
\begin{blocLivre}{cor-11-1}
\begin{corrige}{1}{Calculer une proportion}
\demarche{Diviser la partie par le total et changer d’écriture.}
La référence est l’ensemble des 240 élèves : $84/240=\res{7/20}=\res{0{,}35}=\res{35\,\%}$. Les trois écritures décrivent le même rapport. L’effectif 84 reste un nombre d’élèves, tandis que 0,35 est une proportion sans unité.
\end{corrige}
\end{blocLivre}
\begin{blocLivre}{cor-11-2}
\begin{corrige}{2}{Appliquer un pourcentage}
\demarche{Calculer la part retirée ou utiliser directement le coefficient restant.}
La remise vaut $80\times0{,}15=12$ euros. Le prix final est $80-12=\res{68\text{ euros}}$, ou directement $80\times0{,}85$. Le coefficient multiplicateur est 0,85 car il reste 85\,\% du prix initial.
\end{corrige}
\end{blocLivre}
\begin{blocLivre}{cor-11-3}
\begin{corrige}{3}{Retrouver le total}
\demarche{Diviser l’effectif connu par sa proportion.}
Notons $N$ le total. On a $0{,}18N=54$, donc $N=54/0{,}18=\res{300}$. Contrôle : $18\,\%$ de 300 vaut bien 54. Multiplier 54 par 18\,\% aurait calculé une partie d’une partie et non le total cherché.
\end{corrige}
\end{blocLivre}
\begin{blocLivre}{cor-11-4}
\begin{corrige}{4}{Du départ à l’arrivée}
\demarche{Utiliser la valeur de départ comme dénominateur du taux.}
La variation absolue vaut $150-120=\res{30}$. Le coefficient est $150/120=\res{1{,}25}$. Le taux est $1{,}25-1=\res{25\,\%}$. Ces trois réponses sont complémentaires : une différence, un rapport et une variation relative.
\end{corrige}
\end{blocLivre}
\begin{blocLivre}{cor-11-5}
\begin{corrige}{5}{Une baisse}
\demarche{Conserver le signe de la variation et la bonne référence.}
La variation vaut $210-250=-40$ kWh. Rapportée à la valeur initiale, elle donne $-40/250=-0{,}16$, soit $\res{-16\,\%}$. Le coefficient multiplicateur est 0,84. Le signe négatif exprime la baisse, et la valeur absolue 16\,\% sa proportion.
\end{corrige}
\end{blocLivre}
\begin{blocLivre}{cor-11-6}
\begin{corrige}{6}{Un pourcentage de pourcentage}
\demarche{Identifier les ensembles emboîtés avant de multiplier les proportions.}
Les élèves de l’option sont $800\times0{,}45=360$. Ceux de théâtre sont $360\times0{,}40=\res{144}$. La proportion dans tout le lycée vaut $0{,}45\times0{,}40=0{,}18$, soit $\res{18\,\%}$. Les 40\,\% de l’énoncé portaient seulement sur les 360 élèves de l’option.
\label{fig-cor-11-6}
\figureGeo[corrige]{proportions-emboitees}
\end{corrige}
\end{blocLivre}
\begin{blocLivre}{cor-11-7}
\begin{corrige}{7}{Deux évolutions successives}
\demarche{Multiplier les coefficients plutôt qu’additionner les taux.}
Les coefficients sont 1,10 puis 0,85. Leur produit vaut $0{,}935$. Le prix final est $200\times0{,}935=\res{187\text{ euros}}$ ; le taux global est $0{,}935-1=\res{-6{,}5\,\%}$. Additionner les deux taux aurait ignoré que la seconde évolution porte sur un prix déjà modifié.
\end{corrige}
\end{blocLivre}
\begin{blocLivre}{cor-11-8}
\begin{corrige}{8}{L’évolution réciproque}
\demarche{Inverser le coefficient multiplicateur.}
Le coefficient de hausse est 1,25. Le coefficient réciproque vaut $1/1{,}25=0{,}8$ ; le taux correspondant est donc $\res{-20\,\%}$. Contrôle sur 100 individus : la hausse donne 125, puis une baisse de 20\,\% de 125 retire 25 et ramène à 100.
\end{corrige}
\end{blocLivre}
\begin{blocLivre}{cor-11-9}
\begin{corrige}{9}{Points et pourcentage}
\demarche{Distinguer différence de taux et évolution relative de ce taux.}
La différence est $72-60=\res{12\text{ points}}$. La hausse relative vaut $(72-60)/60=0{,}20$, soit $\res{20\,\%}$. Les deux formulations sont correctes si leur sens est précisé. Dire seulement « une hausse de 12\,\% » serait ici ambigu et ne correspondrait pas au taux relatif.
\end{corrige}
\end{blocLivre}
\begin{blocLivre}{cor-11-10}
\begin{corrige}{10}{Une proportion globale}
\demarche{Revenir aux effectifs avant de calculer le pourcentage global.}
Les effectifs de réussite sont 16 et 18. Au total, il y a 34 réussites sur 50 élèves : $\res{68\,\%}$. La moyenne simple $(80+60)/2=70\,\%$ serait incorrecte, car les classes n’ont pas le même effectif. On peut aussi utiliser une moyenne pondérée par 20 et 30.
\end{corrige}
\end{blocLivre}
\begin{blocLivre}{cor-11-11}
\begin{corrige}{11}{Retrouver un prix hors taxe}
\demarche{Retrouver le départ en divisant par le coefficient.}
Le TTC vaut $1{,}20$ fois le hors taxe. Donc $HT=144/1{,}20=\res{120\text{ euros}}$. La TVA vaut $144-120=\res{24\text{ euros}}$. Retirer 20\,\% de 144 donnerait 115,20 euros et utiliserait la mauvaise référence : la taxe est calculée sur le HT.
\end{corrige}
\end{blocLivre}
\begin{blocLivre}{cor-11-12}
\begin{corrige}{12}{Une baisse suivie d’un retour}
\demarche{Calculer un taux réciproque et préciser son écriture en pourcentage.}
Après la perte, le coefficient est 0,70. Le retour exige $1/0{,}70=10/7$. Le taux de hausse est donc $10/7-1=\res{3/7}$, soit environ $\res{42{,}9\,\%}$. La hausse est supérieure à 30\,\% parce qu’elle part d’une valeur devenue plus petite.
\end{corrige}
\end{blocLivre}
\begin{blocLivre}{cor-11-13}
\begin{corrige}{13}{Un seuil après plusieurs hausses}
\demarche{Itérer le coefficient jusqu’au premier franchissement.}
Les valeurs aux années 1 à 4 sont 550 ; 605 ; 665,50 ; 732,05 euros. Le premier dépassement est donc à $\res{4\text{ ans}}$. Le coefficient 1,1 est supérieur à 1 : les valeurs augmentent, ce qui justifie que les années précédentes ne pouvaient pas dépasser 700 après l’échec à l’année 3.
\end{corrige}
\end{blocLivre}
\begin{blocLivre}{cor-11-14}
\begin{corrige}{14}{Des données de population}
\demarche{Lire le périmètre statistique et choisir explicitement la référence.}
L’écart est $327734-267991=\res{59743\text{ habitants}}$. La référence demandée est Bordeaux : le rapport d’écart vaut $59743/267991\approx0{,}22293$. Nantes compte donc environ $\res{22{,}3\,\%}$ d’habitants de plus dans ce périmètre et à cette date. Inverser la comparaison donnerait un autre pourcentage.
\end{corrige}
\end{blocLivre}
\begin{blocLivre}{cor-11-15}
\begin{corrige}{15}{Un taux constant à retrouver}
\demarche{Traduire la répétition du même coefficient par un carré.}
Notons $C>1$ le coefficient de chaque hausse. Le coefficient global vaut $C^2=1{,}21$. Comme $C>0$, $C=\sqrt{1{,}21}=1{,}1$. Chaque hausse est donc de $\res{10\,\%}$. La valeur $-1{,}1$, autre solution algébrique du carré, ne représente pas un coefficient de hausse positif.
\end{corrige}
\end{blocLivre}
